\paragraph{Attention-based inter-agent communication.} CommNet
\citep{sukhbaatar2016commnet} broadcasts mean-pooled hidden states; TarMAC
\citep{das2019tarmac} introduces signature-based targeting so that receivers
weight senders; IC3Net \citep{singh2018ic3net} gates whether an agent speaks at
all; ATOC \citep{jiang2018atoc} forms bounded communication groups; DGN
\citep{jiang2018dgn} restricts attention to a spatial neighbourhood; MAT
\citep{wen2022mat} casts the joint policy as a sequence model. These works
motivate restricted communication with a cost argument --- bandwidth,
complexity, or ``receiving a large amount of information $\ldots$ incurs high
computational complexity'' \citep{jiang2018dgn}.

\paragraph{The audit.} We checked whether that argument is ever accompanied by
a measurement. For each paper we read the full text and recorded (i) whether it
makes a computational-efficiency or complexity claim about its mechanism and
(ii) whether it reports any measured temporal quantity --- time, speed,
throughput, or achieved FLOP/s --- attributable to that mechanism. A statement
such as ``training took a few days'' does not qualify and is recorded as not
qualifying. Table~\ref{tab:audit} reports the result.

\begin{table}[t]
\centering
\small
\caption{Does the paper measure what it claims to save? ``Reports latency''
requires a measured temporal quantity attributable to the mechanism. The
comparison group is the efficient-attention and sparse mixture-of-experts
literature from which the MARL complexity argument is borrowed.}
\label{tab:audit}
\begin{tabular}{llcc}
\toprule
\textbf{Paper} & \textbf{Efficiency claim} & \textbf{Type} & \textbf{Reports latency} \\
\midrule
\multicolumn{4}{l}{\emph{Multi-agent communication / attention}} \\
CommNet \citep{sukhbaatar2016commnet}      & none explicit          & ---          & no \\
TarMAC \citep{das2019tarmac}               & compact messages       & qualitative  & no \\
IC3Net \citep{singh2018ic3net}             & fewer messages         & counts       & no \\
ATOC \citep{jiang2018atoc}                 & bounded group size     & counts       & no \\
DGN \citep{jiang2018dgn}                   & neighbourhood scaling  & asymptotic   & no \\
MAT \citep{wen2022mat}                     & linear in $N$          & asymptotic   & \textbf{yes} \\
\midrule
\multicolumn{4}{l}{\emph{Efficient attention / sparse mixture-of-experts}} \\
Reformer \citep{kitaev2020reformer}        & $\mathcal{O}(L\log L)$ & asymptotic   & \textbf{yes} \\
Longformer \citep{beltagy2020longformer}   & $\mathcal{O}(L)$       & asymptotic   & \textbf{yes} \\
BigBird \citep{zaheer2020bigbird}          & $\mathcal{O}(L)$       & asymptotic   & no \\
Performer \citep{choromanski2021performers}& $\mathcal{O}(L)$       & asymptotic   & \textbf{yes} \\
Switch \citep{fedus2022switch}             & FLOPs-matched          & FLOPs        & \textbf{yes} \\
Sparse MoE \citep{shazeer2017moe}          & large capacity, low cost & FLOPs      & \textbf{yes} \\
\bottomrule
\end{tabular}
\end{table}

One of six multi-agent communication papers reports a measured temporal
quantity for its mechanism, against five of six in the literature the argument
is imported from. We stress what this table is and is not. It is an audit of
\emph{reporting practice}, not of correctness: a ``no'' means the efficiency
argument rests on operation counts or asymptotics alone, not that the argument
is wrong. It also has an obvious selection caveat --- these are the papers this
line of work cites, not a random sample of MARL.

The asymmetry is nonetheless informative, because the borrowed argument does
not transfer cleanly. The efficient-attention results hold at sequence lengths
of $10^3$--$10^5$, where attention is genuinely the bottleneck; MARL
communication operates at $N \lesssim 10^2$, where it is not. And the borrowed
methods are structurally sparse, so they escape the ceiling of
Section~\ref{sec:analysis} that data-dependent top-$k$ is subject to.

\paragraph{Measurement discipline in the efficient-attention literature.} One
precedent deserves particular mention because it anticipates the distinction
this paper turns on. Longformer \citep{beltagy2020longformer} reports runtime
and memory for \emph{three} implementations of the same attention pattern: a
loop version, a vectorised ``chunks'' version, and a custom CUDA kernel. The
authors note that the chunked variant consumes roughly twice the optimal memory
because it ``computes some zero values'' --- that is, they measured a masked
implementation separately from a genuinely sparse one, and reported the
difference. FlashAttention \citep{dao2022flashattention} and memory-efficient
attention \citep{rabe2021selfattention} make the same point from the other
direction: the winning implementations of dense attention reduce memory traffic
without changing the operation count at all, which is precisely the axis a
FLOPs-based argument cannot see.

\paragraph{Sparse and budgeted attention mechanisms.} Sparsemax
\citep{martins2016sparsemax} and $\alpha$-entmax \citep{peters2019entmax}
obtain exact zeros from the normalisation itself rather than a hard cutoff;
Gumbel-softmax \citep{jang2017gumbel} provides the relaxation used by learned
discrete gates; constrained-RL formulations
\citep{achiam2017cpo,stooke2020pidlagrangian} supply the dual machinery for
enforcing a communication budget rather than trading it against return at a
fixed weight. We implement entmax and a Lagrangian budget as arms in
Section~\ref{sec:training}.

\paragraph{Empirical rigour in MARL.} Our training protocol follows the
recommendations of \citet{papoudakis2021epymarl} for MARL benchmarking, and we
report effect sizes and intervals rather than bare means. Our environments come
from the PettingZoo/MPE family \citep{terry2020pettingzoo}, and our learner is
MAPPO \citep{yu2022mappo}.
